\chapter{Trigonométrie : cercle, radians et angles associés}\label{C10}
\begin{repere}{Le parcours du chapitre}
\textbf{Objectifs.} Passer des degrés aux radians ; placer des points ; calculer sinus et cosinus remarquables ; exploiter les symétries du cercle.

\textbf{Prérequis.} Cercle, Pythagore, cosinus et sinus dans un triangle rectangle, symétries, fractions.

\textbf{Repère annuel.} 2–11 mars : 8 h\end{repere}
\section*{Diagnostic : six questions, 10 minutes}\label{C10-D}
Répondre sans consulter le cours. Indiquer les questions pour lesquelles une aide est nécessaire. Ce diagnostic ne donne pas lieu à une note.
\begin{enumerate}
\item Périmètre d’un cercle de rayon 1 ?
\item Le quart d’un tour vaut combien de degrés ?
\item Calculer $\sqrt {\dfrac{1}{2}}$.
\item Abscisse et ordonnée du symétrique de $(a;b)$ par rapport à l’axe des abscisses ?
\item Simplifier $\dfrac{5\pi}{3}-2\pi $.
\item Dans un triangle rectangle, côté opposé 3 et hypoténuse 5 : sinus ?
\end{enumerate}
\paragraph{Après le diagnostic.} Reprendre les notions signalées avec le professeur ; refaire une question analogue, puis expliquer la règle utilisée.
\paragraph{Oral et logique.} Présenter une égalité d’angles associés en nommant la symétrie qui la justifie ; expliquer ce qu’un radian mesure.
\clearpage
\section*{Chercher avant de formaliser}\label{C10-RECH}
Construire, sans rapporteur, un point d’angle $\dfrac{\pi}{3}$ puis un point d’angle $\dfrac{\pi}{4}$. Justifier les valeurs exactes de leurs coordonnées à partir des figures.
\paragraph{Organisation.} Recherche individuelle 5 min, confrontation à deux 5 min, mise en commun 5 min. Garder une trace d’un essai, même non abouti. Une conjecture doit ensuite être justifiée.
\section{Mesurer un angle en radians}\label{C10-U1}
\begin{repere}{Définition et résultats}
Le cercle trigonométrique est le cercle de centre O et de rayon 1, orienté dans le sens inverse des aiguilles d’une montre. Un radian correspond à un arc de longueur 1 sur ce cercle. Un tour vaut $2\pi$ radians, un demi-tour $\pi$. Sur un cercle de rayon r, une longueur d’arc vaut $r\theta$ pour un angle positif $\theta$ exprimé en radians.
\end{repere}
\paragraph{Statut et justification.} Définitions et proportionnalité entre angle au centre et longueur d’arc. La conversion est $180^\circ=\pi$ radians ; on multiplie les degrés par $\dfrac{\pi}{180}$.
\exemple{$60^\circ =\dfrac{\pi}{3}$ ; $\dfrac{3\pi}{2}=270^\circ $. Un angle de $\dfrac{\pi}{3}$ sur un cercle de rayon 6 cm donne un arc de $2\pi$ cm.}\label{C10-EX1}
\begin{erreur}La longueur d’un arc n’est pas celle de la corde. En calcul trigonométrique, vérifier le mode degré ou radian de l’outil.\end{erreur}
\application{Convertir $45^\circ$, $120^\circ$ ; donner l’arc de rayon 2 m et angle $\dfrac{3\pi}{4}$.}\label{C10-AP1}
\section{Enroulement et image d’un réel}\label{C10-U2}
\begin{repere}{Définition et résultats}
À chaque réel x correspond un point M du cercle obtenu par un parcours orienté de longueur x depuis I(1;0). Deux réels ont la même image si leur différence est un multiple entier de $2\pi$. Un nombre négatif indique un parcours dans le sens horaire.
\end{repere}
\paragraph{Statut et justification.} Ajouter ou retirer un nombre entier de tours conserve le point d’arrivée. Réciproquement, revenir au même point exige un nombre entier de tours.
\exemple{$\dfrac{13\pi}{3}=4\pi +\dfrac{\pi}{3}$ a la même image que $\dfrac{\pi}{3}$. $-\dfrac{\pi}{2}$ a la même image que $\dfrac{3\pi}{2}$ : point $(0;-1)$.}\label{C10-EX2}
\begin{erreur}Ajouter $\pi$ donne le point opposé, pas le même point.\end{erreur}
\application{Placer les images de $\dfrac{7\pi}{2}$ et $-\dfrac{5\pi}{3}$ en les ramenant à un tour.}\label{C10-AP2}
\section{Sinus, cosinus et valeurs remarquables}\label{C10-U3}
\begin{repere}{Définition et résultats}
Si M est l’image de x, $M=(\cos x;\sin x)$. Donc $-1\leq\cos x,\sin x\leq1$ et $\cos^2x+\sin^2x=1$. Pour $x=0,\dfrac{\pi}{6},\dfrac{\pi}{4},\dfrac{\pi}{3},\dfrac{\pi}{2}$, les cosinus sont $1,\dfrac{\sqrt 3}{2},\dfrac{\sqrt 2}{2},\dfrac{1}{2},0$ et les sinus $0,\dfrac{1}{2},\dfrac{\sqrt 2}{2},\dfrac{\sqrt 3}{2},1$.
\end{repere}
\paragraph{Statut et justification.} À $\dfrac{\pi}{4}$, les deux coordonnées sont égales et positives ; Pythagore donne $2a^2=1$, donc $a=\dfrac{\sqrt 2}{2}$. Un triangle équilatéral de côté 1, coupé par sa hauteur, a demi-base $\dfrac{1}{2}$ et hauteur $\dfrac{\sqrt 3}{2}$ : $\cos (\dfrac{\pi}{3})=\dfrac{1}{2}$, $\sin (\dfrac{\pi}{3})=\dfrac{\sqrt 3}{2}$. L’angle complémentaire donne $\dfrac{\pi}{6}$.
\exemple{Le point d’angle $\dfrac{\pi}{3}$ est $(\dfrac{1}{2};\dfrac{\sqrt 3}{2})$. Le produit scalaire de deux vecteurs unitaires faisant cet angle vaut son cosinus, donc $\dfrac{1}{2}$. Les chapitres de géométrie se raccordent ainsi au cercle.}\label{C10-EX3}
\par\medskip\noindent\begin{minipage}{\linewidth}\centering
\begin{tikzpicture}[scale=2.8,>=Stealth]\draw[->](-1.3,0)--(1.45,0)node[right]{$\cos x$};\draw[->](0,-1.25)--(0,1.35)node[above]{$\sin x$};\draw[bleu,thick](0,0)circle(1);\draw[vert,very thick](0,0)--(.5,.8660254);\draw[dashed](.5,0)--(.5,.8660254)--(0,.8660254);\draw[->](.3,0)arc(0:60:.3);\node[fill=white,inner sep=1pt]at(.46,.22){$\dfrac\pi3$};\fill(.5,.8660254)circle(.7pt)node[above right]{$M$};\node[below,fill=white,inner sep=2pt]at(.5,0){$\dfrac12$};\node[left,fill=white,inner sep=2pt]at(0,.8660254){$\dfrac{\sqrt3}{2}$};\node[below left]at(0,0){$O$};\end{tikzpicture}
\par\medskip Pour l’angle $\dfrac\pi3$, les projections de M donnent directement son cosinus et son sinus. Les fractions sont placées hors des traits du dessin.\end{minipage}\par\medskip

\begin{erreur}Le cosinus est l’abscisse, le sinus l’ordonnée. L’identité donne une valeur absolue : le quadrant détermine le signe.\end{erreur}
\application{Si $\cos x=\dfrac{3}{5}$ et l’image de x est dans le quatrième quadrant, donner $\sin x$.}\label{C10-AP3}
\section{Angles associés et symétries}\label{C10-U4}
\begin{repere}{Définition et résultats}
Les symétries du cercle donnent $\cos(-x)=\cos x$, $\sin(-x)=-\sin x$ ; $\cos(\pi-x)=-\cos x$, $\sin(\pi-x)=\sin x$ ; $\cos(\pi+x)=-\cos x$, $\sin(\pi+x)=-\sin x$. Ajouter $2k\pi$, k entier, ne change ni sinus ni cosinus.
\end{repere}
\paragraph{Statut et justification.} Chaque égalité se lit sur les coordonnées des points symétriques : axe horizontal pour $-x$, axe vertical pour $\pi-x$, centre O pour $\pi+x$. Il ne s’agit pas d’un catalogue sans justification.
\exemple{$\cos (\dfrac{5\pi}{6})=-\dfrac{\sqrt 3}{2}$, $\sin (\dfrac{5\pi}{6})=\dfrac{1}{2}$ par symétrie de $\dfrac{\pi}{6}$. $\cos (-\dfrac{\pi}{4})=\dfrac{\sqrt 2}{2}$, $\sin (-\dfrac{\pi}{4})=-\dfrac{\sqrt 2}{2}$.}\label{C10-EX4}
\begin{erreur}Deux angles de même cosinus n’ont pas nécessairement le même sinus. Réduire l’angle puis lire les signes.\end{erreur}
\application{Calculer les deux coordonnées des images de $\dfrac{4\pi}{3}$ et $\dfrac{7\pi}{4}$.}\label{C10-AP4}
\section{Approximer pi et réinvestir}\label{C10-U5}
\begin{repere}{Définition et résultats}
Un polygone régulier inscrit à n côtés dans le cercle unité a demi-périmètre $n\sin (\dfrac{\pi}{n})$ ; un polygone circonscrit a demi-périmètre $n\tan (\dfrac{\pi}{n})$, avec $\tan \theta =\dfrac{\sin \theta}{\cos}\theta $ ici. La méthode d’Archimède double le nombre de côtés et calcule les nouvelles longueurs par Pythagore, sans connaître pi à l’avance.
\end{repere}
\paragraph{Statut et justification.} Le périmètre inscrit est inférieur à celui du cercle et le circonscrit supérieur : encadrement géométrique admis. Si c est la corde d’un côté, poser $h=\sqrt {1-\dfrac{c^2}{4}}$ ; la corde du demi-angle vaut $\dfrac{c}{\sqrt {2(1+h)}}$. Cette formule vient de Pythagore et évite une soustraction de réels proches.
\exemple{Pour l’hexagone inscrit, c=1 et demi-périmètre 3. Le circonscrit donne $2\sqrt3\approx3{,}4641$. En doublant les côtés, on resserre l’encadrement. Le programme fourni part de n=6 et c=1.}\label{C10-EX5}
\begin{erreur}Utiliser directement une fonction sinus avec pi déjà connu vérifierait un calcul, mais ne construirait pas une approximation indépendante de pi.\end{erreur}
\application{Pour n=6 et c=1, calculer la nouvelle corde après doublement.}\label{C10-AP5}
\section*{Parcours d’exercices}\label{C10-PARC}
Commencer par les exercices 1 et 2. Choisir ensuite un exercice 3 ou 4 sans aide. Les exercices 5 à 7 sont une réserve de transfert, de recherche ou de programmation : le professeur en désigne un. L’exercice 8 se reprend à distance.
\begingroup
\exercice{C10-E1}{Convertir et mesurer}
On utilise le cercle trigonométrique pour mesurer les angles en radians.
\begin{enumerate}[label=\arabic*.]
\item Convertir $150^\circ$ et $-90^\circ$ en radians, en faisant apparaître le facteur de conversion.
\item Un secteur circulaire a un rayon de 4 cm et un angle de $\dfrac\pi3$ radians. Calculer la longueur exacte de son arc.
\item Expliquer pourquoi une mesure d’angle peut être négative alors qu’une longueur d’arc géométrique est positive.
\end{enumerate}
\exercice{C10-E2}{Réduire les angles}
On considère les réels $\dfrac{17\pi}{6}$, $-\dfrac{7\pi}{4}$ et $5\pi$.
\begin{enumerate}[label=\arabic*.]
\item Pour chacun, déterminer un représentant dans $[0;2\pi[$ ayant la même image sur le cercle trigonométrique. Faire apparaître le nombre de tours ajoutés ou retirés.
\item Placer les trois images sur un cercle et les identifier par les réels de départ.
\end{enumerate}
\exercice{C10-E3}{Valeurs exactes}
On cherche les valeurs exactes du cosinus et du sinus de $\dfrac{2\pi}{3}$, $-\dfrac\pi6$ et $\dfrac{5\pi}{4}$.
\begin{enumerate}[label=\arabic*.]
\item Placer les trois images sur le cercle trigonométrique. Pour chacune, préciser le signe attendu de ses deux coordonnées.
\item Calculer les valeurs exactes en utilisant un angle remarquable et une symétrie. Présenter la réponse dans un tableau : angle, cosinus, sinus.
\end{enumerate}
\exercice{C10-E4}{Retrouver le signe}
Un réel $x$ vérifie $\sin x=\dfrac5{13}$. Son image sur le cercle trigonométrique se situe dans le deuxième quadrant, donc à gauche de l’axe des ordonnées et au-dessus de l’axe des abscisses.
\begin{enumerate}[label=\arabic*.]
\item Déterminer la valeur exacte de $\cos x$. Justifier à la fois sa valeur absolue et son signe.
\item Expliquer pourquoi la seule donnée de $\sin x$ n’aurait pas suffi à choisir le signe de $\cos x$.
\end{enumerate}
\exercice{C10-E5}{Une roue}
Une roue de rayon $0{,}30$ m roule sur un sol horizontal sans glisser. Elle effectue une rotation d’angle $\dfrac{5\pi}{3}$ radians. On cherche la distance parcourue par son axe dans la direction du déplacement.
\begin{enumerate}[label=\arabic*.]
\item Calculer cette distance sous forme exacte en expliquant le lien avec la longueur d’arc déroulée.
\item Expliquer pourquoi on ne doit pas utiliser la longueur d’une corde du cercle. La trajectoire d’un point fixé sur le bord de la roue, vue depuis le sol, est-elle celle de l’axe ?
\end{enumerate}
\exercice{C10-E6}{Corde et produit scalaire}
Sur le cercle trigonométrique de centre $O$, le point $A$ est l’image de 0 et le point $B$ celle de $\dfrac{2\pi}{3}$.
\begin{enumerate}[label=\arabic*.]
\item Déterminer les coordonnées exactes de $A$ et $B$, puis calculer la longueur $AB$.
\item Retrouver cette longueur en appliquant la relation d’Al-Kashi au triangle $OAB$. Préciser les longueurs et l’angle utilisés.
\end{enumerate}
\exercice{C10-E7}{Lire une équation sur le cercle}
On cherche les réels $x$ de l’intervalle $[0;2\pi[$ dont l’image sur le cercle trigonométrique a pour abscisse $\dfrac12$.
\begin{enumerate}[label=\arabic*.]
\item Traduire la condition sur l’abscisse par une équation faisant intervenir $\cos x$.
\item Tracer la droite verticale correspondante sur un cercle, repérer ses intersections et donner les deux solutions. Vérifier qu’elles appartiennent à l’intervalle demandé.
\end{enumerate}
\exercice{C10-E8}{Reprise différée}
On considère l’angle de mesure $\dfrac{25\pi}{4}$ radians. Les réponses doivent être exactes, sans approximation décimale.
\begin{enumerate}[label=\arabic*.]
\item Réduire cet angle en faisant apparaître un nombre entier de tours complets.
\item En déduire son cosinus et son sinus. Justifier les valeurs à l’aide du cercle trigonométrique.
\end{enumerate}
\paragraph{Travail personnel.} Effectuer uniquement la tâche désignée à l’issue de la séance, dans la progression datée : 10–20 min maximum. Les exercices dépendent des notions déjà étudiées ; certaines questions sont donc différées. Reprendre l’exercice 8 une semaine après le chapitre (10 min). Si le temps prévu est dépassé, noter la difficulté et apporter les essais ; ne pas accumuler toute la réserve.
\endgroup
\Needspace{12\baselineskip}\section*{À faire dans le manuel CQFD}
Ces exercices se traitent dans votre manuel. La progression indique la séance ; ils remplacent le travail personnel prévu ce jour-là (10 à 20 minutes au total).
\paragraph{CQFD, exercice 32 p. 220, questions a et b.}\label{C10-M1} Utiliser la valeur exacte fournie et les angles associés, sans calculatrice.
\paragraph{CQFD, exercice 39 p. 220, questions 1, 2, 3 et 5.}\label{C10-M2} Traiter les quatre questions de calcul et de lecture du cercle ; justifier les réponses sur la copie.
\begingroup\small
\clearpage\section*{QCM 1 : comprendre et choisir}\label{C10-Q1}
Cinq questions pédagogiques à réponse unique. Choisir une réponse par question ; expliquer ensuite un distracteur. Ce format de chapitre ne reproduit pas le format officiel de l’EAM.
\begin{enumerate}
\item $180^\circ$ vaut
\begin{enumerate}[label=\Alph*.,itemsep=2pt]
\item $2\pi$
\item $\pi$
\item $\dfrac{\pi}{2}$
\item $1$
\end{enumerate}
\item $\cos (\dfrac{\pi}{3})=$
\begin{enumerate}[label=\Alph*.,itemsep=2pt]
\item $\dfrac{1}{2}$
\item $\dfrac{\sqrt 3}{2}$
\item $0$
\item $1$
\end{enumerate}
\item $\sin (-\dfrac{\pi}{2})=$
\begin{enumerate}[label=\Alph*.,itemsep=2pt]
\item $1$
\item $0$
\item $\dfrac{1}{2}$
\item $-1$
\end{enumerate}
\item $\cos^2x+\sin^2x=$
\begin{enumerate}[label=\Alph*.,itemsep=2pt]
\item $0$
\item $2$
\item $1$
\item $x$
\end{enumerate}
\item Un même point correspond à des angles différant de
\begin{enumerate}[label=\Alph*.,itemsep=2pt]
\item $\pi$
\item $2k\pi$, k entier
\item $\dfrac{\pi}{2}$
\item $1$
\end{enumerate}
\end{enumerate}
\endgroup
\begingroup\small
\Needspace{18\baselineskip}\section*{QCM 2 : comprendre et choisir}\label{C10-Q2}
Cinq questions pédagogiques à réponse unique. Choisir une réponse par question ; expliquer ensuite un distracteur. Ce format de chapitre ne reproduit pas le format officiel de l’EAM.
\begin{enumerate}
\item $\cos(\pi-x)=$
\begin{enumerate}[label=\Alph*.,itemsep=2pt]
\item $\cos x$
\item $\sin x$
\item $-\cos x$
\item $-\sin x$
\end{enumerate}
\item $\sin (\dfrac{\pi}{4})=$
\begin{enumerate}[label=\Alph*.,itemsep=2pt]
\item $\dfrac{\sqrt 2}{2}$
\item $\dfrac{1}{2}$
\item $\dfrac{\sqrt 3}{2}$
\item $\sqrt2$
\end{enumerate}
\item Image de $\dfrac{3\pi}{2}$ :
\begin{enumerate}[label=\Alph*.,itemsep=2pt]
\item $(1;0)$
\item $(-1;0)$
\item $(0;1)$
\item $(0;-1)$
\end{enumerate}
\item Arc de rayon 3 et angle 2 radians : longueur
\begin{enumerate}[label=\Alph*.,itemsep=2pt]
\item $\dfrac{2}{3}$
\item $6$
\item $\dfrac{3}{2}$
\item $6\pi$
\end{enumerate}
\item Intérieur du deuxième quadrant : signe du cosinus
\begin{enumerate}[label=\Alph*.,itemsep=2pt]
\item positif
\item nul partout
\item négatif
\item indéterminable
\end{enumerate}
\end{enumerate}
\subsection*{Vrai ou faux ? Justifier.}\label{C10-VF}
\begin{enumerate}
\item $\cos(-x)=-\cos x$ pour tout x.
\item $\sin(x+2\pi)=\sin x$.
\item La donnée de $\sin x$ détermine toujours $\cos x$.
\end{enumerate}
\begin{repere}{Bilan personnel}Je réduis les angles par tours complets ; je lis cosinus et sinus dans le bon ordre ; je justifie les signes ; je conserve les valeurs exactes.

Pour chaque objectif, indiquer : acquis sans aide ; réussi avec aide ; à retravailler. Écrire un exemple et une erreur à éviter.\end{repere}
\endgroup
\ifprof\corrsection{10}\label{C10-CORR}
\textbf{Usage professeur.} Les réponses qui suivent concernent les tâches non résolues dans le cours. Faire d’abord expliciter la démarche ; fournir une aide graduée avant la solution. Les durées sont des estimations à étalonner en classe.
\subsection*{Diagnostic et remédiation}
\begin{enumerate}
\item $2\pi$. \textit{Action :} Relier rayon, diamètre et périmètre.
\item $90^\circ$. \textit{Action :} Partager un cercle en quatre secteurs.
\item $\dfrac{\sqrt 2}{2}$. \textit{Action :} Rationaliser et vérifier le carré.
\item $(a;-b)$. \textit{Action :} Tracer les deux points.
\item $-\dfrac{\pi}{3}$. \textit{Action :} Mettre au même dénominateur.
\item $\dfrac{3}{5}$. \textit{Action :} Identifier l’angle de référence.
\end{enumerate}
\subsection*{Activité de recherche}
Triangle équilatéral puis triangle rectangle isocèle. Pythagore donne respectivement hauteur $\dfrac{\sqrt 3}{2}$ et coordonnées égales $\dfrac{\sqrt 2}{2}$. Les signes sont positifs dans le premier quadrant.
\subsection*{Applications immédiates}
\begin{enumerate}
\item $\dfrac{\pi}{4}$, $\dfrac{2\pi}{3}$ ; longueur $\dfrac{3\pi}{2}$ m.
\item $\dfrac{7\pi}{2}$ a l’image de $\dfrac{3\pi}{2}$ ; $-\dfrac{5\pi}{3}$ celle de $\dfrac{\pi}{3}$.
\item $\sin ^2x=\dfrac{16}{25}$ ; le sinus est négatif, donc $-\dfrac{4}{5}$.
\item $(-\dfrac{1}{2};-\dfrac{\sqrt 3}{2})$ et $(\dfrac{\sqrt 2}{2};-\dfrac{\sqrt 2}{2})$.
\item $h=\dfrac{\sqrt 3}{2}$ ; nouvelle corde $\dfrac{1}{\sqrt {2+\sqrt3}}=\sqrt {2-\sqrt3}$.
\end{enumerate}
\subsection*{Exercices : démarche et résultats}
\paragraph{Exercice 10.1 -- Convertir et mesurer.} $\dfrac{5\pi}{6}$ et $-\dfrac{\pi}{2}$ ; arc $\dfrac{4\pi}{3}$ cm. Le signe négatif indique une orientation ; une longueur géométrique est positive.
\paragraph{Exercice 10.2 -- Réduire les angles.} $\dfrac{5\pi}{6}$, $\dfrac{\pi}{4}$, $\pi$. Retirer ou ajouter des multiples de $2\pi$.
\paragraph{Exercice 10.3 -- Valeurs exactes.} $(-\dfrac{1}{2};\dfrac{\sqrt 3}{2})$ ; $(\dfrac{\sqrt 3}{2};-\dfrac{1}{2})$ ; $(-\dfrac{\sqrt 2}{2};-\dfrac{\sqrt 2}{2})$, sous forme (cosinus ; sinus).
\paragraph{Exercice 10.4 -- Retrouver le signe.} $\cos ^2x=1-\dfrac{25}{169}=\dfrac{144}{169}$ ; cosinus négatif, donc $-\dfrac{12}{13}$.
\paragraph{Exercice 10.5 -- Une roue.} Avancement $r\theta =0{,}3(\dfrac{5\pi}{3})=\dfrac{\pi}{2}$ m. Cela correspond au déroulement de l’arc ; la corde relie deux positions sur le cercle et n’a pas cette longueur. Le point matériel sur la roue décrit une trajectoire différente par rapport au sol.
\paragraph{Exercice 10.6 -- Corde et produit scalaire.} $A(1;0)$, $B(-\dfrac{1}{2};\dfrac{\sqrt 3}{2})$ ; $AB^2=\dfrac{9}{4}+\dfrac{3}{4}=3$. Al-Kashi : $1+1-2\cos (\dfrac{2\pi}{3})=3$. Donc $AB=\sqrt3$.
\paragraph{Exercice 10.7 -- Lire une équation sur le cercle.} La verticale d’abscisse $\dfrac{1}{2}$ coupe le cercle aux angles $\dfrac{\pi}{3}$ et $\dfrac{5\pi}{3}$. Les deux points ont des sinus opposés.
\paragraph{Exercice 10.8 -- Reprise différée.} $\dfrac{25\pi}{4}=6\pi +\dfrac{\pi}{4}$. Trois tours complets ne changent pas l’image, donc les deux valeurs valent $\dfrac{\sqrt 2}{2}$.
\paragraph{CQFD, exercice 32 p. 220, questions a et b.} a : $\cos(\dfrac{6\pi}5)=\cos(\pi+\dfrac\pi5)=-\dfrac{1+\sqrt5}{4}$. b : $\cos(-\dfrac\pi5)=\cos(\dfrac\pi5)=\dfrac{1+\sqrt5}{4}$, par parité du cosinus.
\paragraph{CQFD, exercice 39 p. 220, questions 1, 2, 3 et 5.} 1 : $\dfrac\pi6+\dfrac\pi3=\dfrac\pi2$. 2 : oui, $\cos(\dfrac\pi3)=\dfrac12=0{,}5$ est décimal. 3 : la mesure principale dans $]-\pi;\pi]$ est $-\dfrac{5\pi}6$. 5 : $\sin(\dfrac{4\pi}3)=-\dfrac{\sqrt3}2$.
\subsection*{QCM : réponses et raisons}
\textbf{QCM 1.}
\begin{enumerate}
\item B. Demi-tour.
\item A. Triangle équilatéral.
\item D. Point inférieur du cercle.
\item C. Équation du cercle unité.
\item B. Tours complets.
\end{enumerate}
\textbf{QCM 2.}
\begin{enumerate}
\item C. Symétrie verticale.
\item A. Triangle rectangle isocèle.
\item D. Trois quarts de tour.
\item B. $r\theta=6$.
\item C. Abscisse négative à l’intérieur du quadrant.
\end{enumerate}
\subsection*{Vrai ou faux}
\begin{enumerate}
\item Faux : en 0, les deux membres seraient 1 et $-1$. Le cosinus est pair.
\item Vrai : un tour complet conserve l’image.
\item Faux : $\dfrac{\pi}{6}$ et $\dfrac{5\pi}{6}$ ont même sinus mais cosinus opposés.
\end{enumerate}
\paragraph{Différenciation.} Pour un blocage technique, proposer la première étape seulement (domaine, facteur ou univers de référence), puis retirer l’aide. Pour un élève autonome, demander une seconde méthode et le traitement d’un cas limite. À l’oral, autoriser une reformulation initiale dans la langue maîtrisée, puis exiger une phrase mathématique précise en français.
\fi
